% Resumen de la solución
\subsection{Resumen de la solución}\label{sec:1_1_resumen_solucion}

\subsubsection{Descripción general de la plataforma (microservicios Spring Boot, frontends Angular y modelo de IA)}
La Plataforma de Quejas de la Defensoría de los Derechos Politécnicos (DDP) es un sistema web que permite a la comunidad del Instituto Politécnico Nacional presentar quejas y darles seguimiento, y al personal de la Defensoría recibirlas, revisarlas, turnarlas y atenderlas. Está construida con una arquitectura de microservicios y se compone de tres tipos de piezas:

\begin{itemize}
  \item \textbf{Microservicios (backend).} Diez servicios independientes desarrollados en Java 21 con Spring Boot. Cada uno resuelve una parte del negocio, expone una API REST y se ejecuta en su propio contenedor de Podman.
  \item \textbf{Frontends (interfaz web).} Cinco aplicaciones Angular, una por cada tipo de usuario. Se compilan a archivos estáticos (HTML, CSS y JavaScript) y cada una se sirve desde su propio contenedor de nginx.
  \item \textbf{Modelo de inteligencia artificial.} El servicio \texttt{antecedentes-service}, también en Java, busca quejas parecidas en el histórico mediante TF-IDF y similitud coseno, y genera un resumen extractivo de cada queja con TextRank. No requiere GPU ni servicios externos.
\end{itemize}

Todos los servicios guardan su información en una sola base de datos PostgreSQL 16 (\texttt{defensoria\_db}).

\begin{table}[H]
\centering
\caption{Microservicios de la plataforma}
\small
\begin{tabular}{p{4.3cm}p{10.2cm}}
\toprule
\textbf{Servicio} & \textbf{Función} \\
\midrule
\texttt{auth-service} & Inicio de sesión, emisión de tokens JWT, activación de cuentas y recuperación de contraseña \\
\texttt{quejas-service} & Registro y consulta de quejas, folios, credenciales y evidencias \\
\texttt{notificaciones-service} & Avisos dentro de la plataforma y envío de correos \\
\texttt{catalogo-service} & Catálogo de dependencias (unidades académicas) del IPN \\
\texttt{admin-service} & Panel de la Defensora: indicadores, personal, configuración y respaldos \\
\texttt{revision-service} & Validación por la recepcionista y turnado de quejas \\
\texttt{chatbot-service} & Preguntas frecuentes del asistente del portal \\
\texttt{denunciado-service} & Registro de la respuesta y evidencias de la persona denunciada \\
\texttt{primer-contacto-service} & Módulo de Primer Contacto (atención inicial y citas) \\
\texttt{subdefensoria-service} & Módulo de la Subdefensoría (atención de quejas turnadas) \\
\texttt{antecedentes-service} & Modelo de IA: búsqueda de antecedentes y resumen de quejas \\
\bottomrule
\end{tabular}
\end{table}

\begin{table}[H]
\centering
\caption{Frontends de la plataforma}
\small
\begin{tabular}{p{4.3cm}p{3.4cm}p{6.5cm}}
\toprule
\textbf{Frontend} & \textbf{Ruta pública} & \textbf{Usuarios} \\
\midrule
Portal (\texttt{defensoria-web}) & \texttt{/} & Quejosos y personas denunciadas \\
Revisión (\texttt{revision-web}) & \texttt{/revision/} & Login unificado del personal y recepcionistas \\
Administración (\texttt{admin-web}) & \texttt{/admin/} & Defensora y administración del sistema \\
Primer Contacto (\texttt{primer-contacto-web}) & \texttt{/primer-contacto/} & Analistas de Primer Contacto \\
Subdefensoría (\texttt{subdefensoria-web}) & \texttt{/subdefensoria/} & Abogados de la Subdefensoría \\
\bottomrule
\end{tabular}
\end{table}

\subsubsection{Objetivo del despliegue en 3 VPS}
El objetivo del despliegue es dejar la plataforma completa en producción y accesible en \url{https://defensoria-escom.ddns.net}, separando sus componentes en tres servidores virtuales, uno por capa:

\begin{itemize}
  \item \textbf{Seguridad.} Solo el servidor frontend queda expuesto a Internet. El backend acepta conexiones únicamente del frontend, y la base de datos únicamente del backend, de modo que los datos de las quejas nunca quedan directamente al alcance de un atacante externo.
  \item \textbf{Aislamiento de recursos.} La base de datos, los microservicios y los portales no compiten por la misma memoria y CPU; si un servicio consume más de lo previsto, no afecta a las otras capas.
  \item \textbf{Mantenimiento independiente.} Cada capa se puede actualizar, reiniciar o respaldar sin detener a las demás.
  \item \textbf{Reproducibilidad.} Todos los componentes se ejecutan en contenedores Podman creados con scripts, por lo que la instalación se puede repetir de la misma forma en servidores nuevos.
\end{itemize}

\begin{table}[H]
\centering
\caption{Servidores del despliegue}
\small
\begin{tabular}{llll}
\toprule
\textbf{Capa} & \textbf{Hostname} & \textbf{IP pública} & \textbf{Contenido} \\
\midrule
Frontend & vmi3633026 & 169.58.62.111 & router-nginx y 5 frontends \\
Backend & vmi3633021 & 156.67.26.73 & 10 microservicios y modelo de IA \\
Base de datos & vmi3633025 & 169.58.62.99 & PostgreSQL 16 (\texttt{defensoria-db}) \\
\bottomrule
\end{tabular}
\end{table}
